\exercisesection{Exercises for § 6}

In the exercises below (with the exception of Exerc. 3), the assumptions and notation are those of §6.

\begin{enumerate}
\def\labelenumi{\arabic{enumi})}
\item
  Assume that W is irreducible. Let c be a Coxeter transformation of W, let \(\Gamma\) be the subgroup of W generated by c, and let \(\Delta\) be the set of unit vectors orthogonal to an element of H. Show that \(\Gamma\) has l orbits in \(\Delta\) , and that each orbit has h elements. (Argue as in the proof of Prop. 33 of Chap. VI, §1, no. 11.)
\item
  Assume that W is irreducible. Let C be a chamber relative to W, \((\mathrm{H}_{1},\ldots,\mathrm{H}_{l})\) its walls, and \(e_{i}\) a non-zero vector orthogonal to \(H_{i}\) . Assume that \(e_{1},\ldots,e_{r}\) (resp. \(e_{r+1},\ldots,e_{l}\) ) are pairwise orthogonal (cf.~no. 2). For all \(u\in Z\) , define \(H_{u}\) and \(s_{u}\) by \(H_{u}=H_{k}\) if \(u\equiv k\pmod{l}\) and \(s_{u}=s_{H_{u}}\) .
\end{enumerate}

\begin{enumerate}
\def\labelenumi{\alph{enumi})}
\item
  Show that the elements of \(\mathfrak{H}\) are the \(s_1 s_2 \ldots s_{u-1} H_u\) for \(u = 1, 2, \ldots, lh/2\) .
\item
  Let \(s' = s_1 \ldots s_r\) and \(s'' = s_{r+1} \ldots s_l\) , so that \(c = s's''\) is the Coxeter transformation associated to the ordered chamber C. Let \(w_0\) be the element of W that transforms C to -C (cf.~§4, Exerc. 2). Show that, if \(h\) is odd,
\end{enumerate}

\[
w _ {0} = \underbrace {s ^ {\prime} s ^ {\prime \prime} s ^ {\prime} \ldots s ^ {\prime \prime} s ^ {\prime}} _ {h \text {terms}} = \underbrace {s ^ {\prime \prime} s ^ {\prime} s ^ {\prime \prime} \ldots s ^ {\prime} s ^ {\prime \prime}} _ {h \text {terms}} = s ^ {\prime \prime} c ^ {(h - 1) / 2} = c ^ {(h - 1) / 2} s ^ {\prime}.
\]

\begin{enumerate}
\def\labelenumi{\alph{enumi})}
\setcounter{enumi}{2}
\tightlist
\item
  Put \(S = \{s_1, \ldots, s_l\}\) . The pair (W,S) is a Coxeter system. If \(w \in W\) , denote by \(l_S(w)\) the length of \(w\) with respect to S (Chap. IV, §1, no. 1). Show that
\end{enumerate}

\[
l _ {\mathrm{S}} (s ^ {\prime}) = r, l _ {\mathrm{S}} (s ^ {\prime \prime}) = l - r, l _ {\mathrm{S}} (c) = l.
\]

Deduce that, with the assumptions in b).

\[
l _ {\mathrm{S}} (w _ {0}) \leqslant r + \frac {h - 1}{2} l \text { and } l _ {\mathrm{S}} (w _ {0}) \leqslant (l - r) + \frac {h - 1}{2} l.
\]

Show on the other hand that \(l_{S}(w_0) = \mathrm{Card}(\mathfrak{H}) = hl / 2\) (use Exerc. 22 of Chap. IV, §1). Deduce that \(r = l / 2\) and that \((s_1, s_2, \ldots, s_{lh / 2})\) is a reduced decomposition of \(w_0\) .

\(d)\) Show that, if \(h\) is even, \((s_1,\ldots ,s_{lh / 2})\) is a reduced decomposition of \(w_{0} = c^{h / 2}\) . (Same method.)

¶ 3) Let K be a commutative ring, E a free K-module with basis \((e_1, \ldots, e_l)\) , and \(f_1, \ldots, f_l\) elements of the dual of E. Put \(a_{ij} = f_i(e_j)\) . If \(1 \leqslant i \leqslant l\) , let \(s_i\) be the pseudo-reflection \(s_{e_i, f_i}\) (\S 2, Exerc. 1). Then

\[
s _ {i} (e _ {j}) = e _ {j} - a _ {i j} e _ {i}.
\]

Put \(c = s_1 \ldots s_l\) , and \(z_i = c(e_i)\) .

\begin{enumerate}
\def\labelenumi{\alph{enumi})}
\tightlist
\item
  For \(1 \leqslant i, k \leqslant l\) , put
\end{enumerate}

\[
y _ {i} ^ {k} = s _ {1} \dots s _ {k} (e _ {i}) \text { and } y _ {i} = s _ {1} \dots s _ {i - 1} (e _ {i}) = y _ {i} ^ {i - 1}.
\]

Thus \(y_{i}^{0} = e_{i}\) and \(y_{i}^{l} = z_{i}\) .

Show that

\[
y _ {i} ^ {k - 1} - y _ {i} ^ {k} = a _ {k i} y _ {k}.
\]

Deduce the formulas

\[
e _ {i} = y _ {i} + \sum_ {k <   i} a _ {k i} y _ {k}
\]

\[
z _ {i} = y _ {i} - \sum_ {k \geqslant i} a _ {k i} y _ {k}.
\]

\begin{enumerate}
\def\labelenumi{\alph{enumi})}
\setcounter{enumi}{1}
\tightlist
\item
  Let C be the matrix of \(c\) with respect to the basis \((e_i)\) . Let \(U = (u_{ij})\) and \(V = (v_{ij})\) be the matrices defined by
\end{enumerate}

\[
u _ {i j} = \left\{ \begin{array}{l l} a _ {i j} & \text { if } i <   j \\ 0 & \text { otherwise, } \end{array} \right. v _ {i j} = \left\{ \begin{array}{l l} 0 & \text { if } i <   j \\ a _ {i j} & \text { otherwise. } \end{array} \right.
\]

The matrix \(\mathrm{I} + \mathrm{U}\) is invertible with determinant 1. Show that

\[
\mathrm{C} = (\mathrm{I} - \mathrm{V}) (\mathrm{I} + \mathrm{U}) ^ {- 1}.
\]

Deduce that

\[
\det (\lambda I - C) = \det ((\lambda - 1) I + V + \lambda U),
\]

in other words

\[
\det (\lambda \mathrm{I} - \mathrm{C}) = \left| \begin{array}{c c c c c} (\lambda - 1) + a _ {1 1} & \lambda a _ {1 2} & \lambda a _ {1 3} & \dots & \lambda a _ {1 l} \\ a _ {2 1} & \lambda (\lambda - 1) + a _ {2 2} & \lambda a _ {2 3} & \dots & \lambda a _ {2 l} \\ a _ {3 1} & a _ {3 2} & (\lambda - 1) + a _ {3 3} & \dots & \lambda a _ {3 l} \\ \vdots & \vdots & \vdots & \ddots & \vdots \\ a _ {l 1} & a _ {l 2} & a _ {l 3} & \dots & (\lambda - 1) + a _ {l l} \end{array} \right|.
\]

\begin{enumerate}
\def\labelenumi{\alph{enumi})}
\setcounter{enumi}{2}
\tightlist
\item
  Let \(\Gamma = (\mathrm{I},\mathrm{S})\) be the graph whose set of vertices is \(\mathrm{I} = [1,l]\) and whose set S of arrows is the set of subsets \(\{i,j\}\) of I with 2 elements such that \(a_{ij}\neq 0\) or \(a_{ji}\neq 0\) . If \(\alpha = \{i,j\}\) belongs to S, denote by \(\sigma_{\alpha}\) the transposition of \(i\) and \(j\) (considered as an element of the symmetric group \(\mathrm{S_I}\) ), and put \(a_{\alpha} = -a_{ij}a_{ji}\) .
\end{enumerate}

Let \(\mathfrak{B}\) be the set of subsets of \(S\) consisting of the arrows whose vertices are disjoint. If \(X \in \mathfrak{B}\) , denote by \(C(X)\) the set of \(i \in I\) that are not vertices of any arrow of \(X\) , and put \(\sigma_X = \prod_{\alpha \in X} \sigma_\alpha\) , \(a_X = \prod_{\alpha \in X} a_\alpha\) .

Let \(\sigma \in S_{I}\) , and let \(d_{\sigma}\) be the term corresponding to \(\sigma\) in the expansion of the determinant of \((\lambda - 1)I + V + \lambda U\) . Show that, if \(\sigma\) is of the form \(\sigma_{X}\) , with \(X \in \mathfrak{B}\) , then

\[
d _ {\sigma} = a _ {X} \lambda^ {\operatorname{Card} (X)} \prod_ {i \in C (X)} (\lambda - 1 + a _ {i i}).
\]

Assume now that \(\Gamma\) is a forest (Chap. IV, Appendix, no. 3). Show that, if \(\sigma \in S_l\) is not of the form \(\sigma_X\) with \(X \in \mathfrak{B}\) , then \(d_{\sigma} = 0\) . Deduce the formula

\[
\det (\lambda - c) = \sum_ {X \in \mathfrak {B}} a _ {X} \lambda^ {\operatorname{Card} (X)} \prod_ {i \in C (X)} (\lambda - 1 + a _ {i i}).
\]

\begin{enumerate}
\def\labelenumi{\alph{enumi})}
\setcounter{enumi}{3}
\tightlist
\item
  Consider the polynomial
\end{enumerate}

\[
P (X) = \left| \begin{array}{c c c c} X & a _ {1 2} & \ldots & a _ {1 l} \\ a _ {2 1} & X & \ldots & a _ {2 l} \\ \vdots & \vdots & \ddots & \vdots \\ a _ {l 1} & a _ {l 2} & \ldots & X \end{array} \right|.
\]

To the assumptions in \(c\) ), we add that \(a_{ii} = 1\) for all \(i\) . Show that in that case

\[
\det (\lambda^ {2} - c) = \lambda^ {l} P (\lambda + \lambda^ {- 1}).
\]

\begin{enumerate}
\def\labelenumi{\arabic{enumi})}
\setcounter{enumi}{3}
\tightlist
\item
  With the assumptions and notation of no. 1, denote by \(n_{ij}\) the order of \(s_{\mathrm{H}_i}s_{\mathrm{H}_j}\) , and put \(a_{ij} = -2\cos \frac{\pi}{n_{ij}}\) . Then \(a_{ii} = 2\) and \(a_{ij} \leqslant 0\) if \(i \neq j\) , cf.~§3. Show \(^{12}\) , by using the method of the preceding exercise, that
\end{enumerate}

\[
\left| \begin{array}{c c c c} X & a _ {1 2} & \ldots & a _ {1 l} \\ a _ {2 1} & X & \ldots & a _ {2 l} \\ \vdots & \vdots & \ddots & \vdots \\ a _ {l 1} & a _ {l 2} & \ldots & X \end{array} \right| = \prod_ {i = 1} ^ {l} (X - 2 \cos (\pi m _ {i} / h)),
\]

where \(m_{1},\ldots ,m_{l}\) are the exponents of W.

